\subsection{Análisis de sensibilidad a los artefactos residuales}
\label{sec:sensibilidad}

Los cuatro sujetos con contaminación ocular residual no determinan la conclusión principal: sin
ellos, el \dsi{} en kappa del resto asciende a $+0.093$ (tabla~\ref{tab:sensibilidad}).

\begin{table}[H]
\centering
\caption{Sujetos con contaminación residual frente al resto (\eegnet{}, \loso{}, media de tres semillas por \textit{fold}).}
\label{tab:sensibilidad}
\begin{tabular}{lccc}
\toprule
 & \textbf{$n$} & \textbf{Exactitud \loso{}} & \textbf{\dsi{} kappa} \\
\midrule
Residuales & 4  & 0.841 & $-0.190$ \\
Resto      & 50 & 0.754 & $+0.093$ \\
\bottomrule
\end{tabular}
\end{table}

Los sujetos marcados por el control ocular tras la limpieza (\texttt{sub-06}, \texttt{sub-29},
\texttt{sub-32} y \texttt{sub-37}) presentan mayor exactitud y \dsi{} negativo, en la dirección que
predice la hipótesis de que la señal ocular residual los vuelve artificialmente fáciles de clasificar.
Con solo cuatro sujetos la diferencia no es concluyente y se reporta como control.

